\scanpage{26}
1.2. Les $C_n$ forment un syst. projectif de $\mathbb Q$-schémas.
\[
\begin{aligned}
C_{n'}&\longrightarrow C_n\deleted{\unsure{'}}\qquad\text{si }n\mid n'\quad\text{soit }n'=nd\\
(X^d,Y^d,\deleted{X^{\unsure{n}}}Z^d)&\longleftarrow X,Y,Z
\end{aligned}\tag*{(4)}
\]
morphismes compatibles avec les
\[
\unsure{\pi|\Pi}_{n'}\longrightarrow\pi_n\tag*{($4'$)}
\]
\margin{\uline{NB} on laisse tomber l'indice $\mathbb Q$}
déduits des hom canoniques
\[
\mu_{n'}\longrightarrow\mu_n\tag*{($4''$)}
\]
Passant à la limite, on trouve
\[
\left\{\begin{aligned}
C_\infty&=\varprojlim_n C_n\\
C_\infty^*&=\varprojlim_n C_n^*
\end{aligned}\right.\tag*{(5)}
\]
sur lequel opère
\[
\begin{aligned}
\pi_\infty=\varprojlim_n\pi_n&\simeq\deleted{T\mathbin{.}}
\overleftarrow{\mu}_\infty\otimes_{\unsure{\mathbb Z|\widehat{\mathbb Z}}}
(\mathbb Z^3/\mathbb Z)\\
&\simeq\overleftarrow{\mu}_\infty\otimes_{\widehat{\mathbb Z}}
(\widehat{\mathbb Z}^3/\widehat{\mathbb Z})
\end{aligned}\tag*{(6)}
\]
où $\overleftarrow{\mu}_\infty=\varprojlim_n\mu_n$.

Ainsi $C_\infty^*$ devient rev. principal de $C_1^*$,
de groupe $\pi_\infty$.
% archive sheet number: 25
